\section{Limitations}\label{sec:limitations}

We collect the main limitations of the framework and of our $D = C \times a_n$ choice in one place so readers can calibrate how strongly to update on the results above.
Several points below recap caveats already made in their natural context and cross-reference those for detail.

\paragraph{$D = C \times a_n$ is a pragmatic choice, not a theoretically derived one.}
The product form was selected because it captures the properties we wanted a diversity score to have (residual surprise remains after the base model has seen previous responses, low-coherence outputs are penalised, and the result stays on a bits-per-byte scale), not because it is derived from an axiomatic information-theoretic setup.
Other scalars respecting the same constraints (weighted integrals of $a_k - a_\infty$, slope-based scores, or coherence applied at a different point on the curve) could equally be defended (see ``The framework admits other metrics'' in Section~\ref{sec:exp-discussion}).
We adopt $D_{Ca_n}$ because it is interpretable and works empirically across our validation regimes, not because it is uniquely correct.
We invite other researchers to build on this work to develop other formulas that characterize the $a_k$ curve in meaningful ways.

\paragraph{Measurement is relative to $\theta$'s perception.}
The metric evaluates diversity \emph{as perceived by the base model $\theta$}: if $\pi$'s outputs differ only in ways $\theta$ cannot distinguish from context, the metric underestimates diversity, and if $\theta$ hallucinates structure, it overestimates redundancy (Section~\ref{sec:motivation}, ``Important caveat'').

\paragraph{Cross-mode learning is model-dependent.}
The off-diagonal sign of the pairwise surprise-reduction matrix changes with $\theta$'s scale (Section~\ref{sec:cross-mode-learning}): $a_k$ curves for the same response set are not directly comparable across different choices of $\theta$, and the $k_0 \propto m$ sigmoid prediction (Appendix~\ref{app:k0-derivation}) only holds when modes appear approximately independent to $\theta$.

\paragraph{$\sigma_\ell$ is a diagnostic, not a diversity signal.}
Coherence spread measures heterogeneity across outputs, not content diversity; it is reported alongside $D$ only to flag mixed-coherence situations (Section~\ref{sec:exp-discussion}).

\paragraph{Metric selection saw the full Tevet benchmark.}
We chose $C \times a_n$ (rather than alternative summaries such as $C \times E$; Appendix~\ref{app:excess-entropy}) after observing its performance on the full Tevet diversity-eval benchmark, without holding out a validation split.
The Tevet numbers in Section~\ref{sec:tevet} may therefore be optimistically biased: the scalar form was selected with visibility into those results.
The OLMo-2-7B post-training experiment (Section~\ref{sec:experiments}) was scored after the metric form was fixed and provides independent validation; the synthetic mode-count experiments (Section~\ref{sec:mode-count-scaling}) informally influenced the choice (had they failed, we would likely have revised the metric) and so are not strictly independent.

\paragraph{External benchmark performance is competitive, not dominant.}
On Tevet and Berant's diversity-eval benchmark, $C \times a_n$ trails SentBERT by \tevetCxAnVsSentBertOCAMinPct\%--\tevetCxAnVsSentBertOCAMaxPct\% OCA across the nine binary tasks (Section~\ref{sec:tevet}); scaling the base model to Qwen3-30B-A3B-Base did not close the gap (Appendix~\ref{app:qwen3-comparison}).
We position $C \times a_n$ as a principled complement to embedding-based metrics, not a replacement.

\paragraph{Prompt format is fixed.}
We use a neutral ``Response A: \ldots, Response B: \ldots'' template throughout and have not attempted prompt optimisation; targeted prompt engineering could plausibly improve discriminative power, especially for larger base models (Section~\ref{sec:exp-discussion}).

\paragraph{Length sensitivity of $D$ in the OLMo RLHF experiment.}
$C$ and $a_n$ are both per-byte quantities, but neither is invariant to response length: in a causal LM, longer responses give $\theta$ more within-response context to predict later tokens, lowering per-byte surprise.
On the OLMo-2-7B post-training experiment (Section~\ref{sec:experiments}) raw response length is not monotone across stages, so the monotone $D$ drop could in principle be driven partly by length rather than diversity per se.
We tested this directly by re-running the experiment with every response truncated to a per-prompt common UTF-8 byte length (the minimum across all 40 responses for that prompt) before re-scoring with the same base $\theta$.
The monotone $D$ drop survives length-matching on both prompt sets, with all three pre-registered H1 contrasts (H1a, H1b, H1c) remaining Bonferroni-significant at $p<0.001$ and effect sizes of the same order of magnitude as the raw analysis.
The single qualification is that $111/300$ prompts had a per-prompt common length below 50 bytes, driven by the base model on AlpacaEval (no instruction-following prior) and by SFT on NB-curated short-answer prompts, and were dropped because length-matching them would compare essentially-empty fragments.
The length-matched effect therefore conditions on prompts where length-matching is well-defined ($150/200$ AlpacaEval, $39/100$ NB-curated).
Full numbers and reproduction commands are in \path{investigations/length_matched_rlhf/VERDICT.md}.

